\chapter{Implantable Cardioverter-Defibrillator Implantation}

\section*{What Is This Test or Procedure?}
An implantable cardioverter-defibrillator (ICD) is a device placed under the skin to detect and treat dangerous ventricular arrhythmias with pacing or an electrical shock.

\section*{Why This Test or Procedure Is Ordered}
It may be recommended after cardiac arrest or for people at high risk of life-threatening ventricular tachycardia or fibrillation, including selected patients with severe cardiomyopathy.

\section*{How This Test or Procedure Is Performed}
Under local anesthesia with sedation, the generator is placed beneath the skin, usually below the collarbone. One or more leads are guided through a vein into the heart and connected to the generator; the system is tested before closure.

\section*{Risks and Follow-up}
Infection, bleeding, pneumothorax, lead displacement, vessel or heart injury, inappropriate shocks, and device malfunction are possible. Regular device checks and battery replacement procedures are required.

\section*{References}
\begin{enumerate}
  \item MedlinePlus. Implantable Cardioverter-Defibrillator. U.S. National Library of Medicine; 2025.
  \item Current Medical Diagnosis and Treatment. McGraw-Hill; 2025.
\end{enumerate}
\clearpage
